% Chapter 5, typeset from lectures/L05: the README is the opener and the closing
% review, and each appendix is a section, in the same order and with the same subsections.

\chapter{Writing a Driver}
\label{c5:ch:driver}\label{c5:ch}

{\large\itshape A driver for a part that does not exist.\par}

\begin{chapterbox}{In this chapter}
\begin{itemize}
  \item The device driver model: \code{struct device}, the API struct, and per-instance config and data.
  \item Defining devices from the devicetree: \code{DT_DRV_COMPAT}, \code{DEVICE_DT_INST_DEFINE()} and
    \code{DT_INST_FOREACH_STATUS_OKAY()}.
  \item Initialization: init levels, priorities, what \code{device_is_ready()} reports after a failure, and
    deferring a device’s initialization with \code{zephyr,deferred-init} and \code{device_init()}.
  \item A binding and a Kconfig option for a new part, and a Kconfig option that turns itself on when the
    devicetree has the part.
  \item Implementing a subsystem API: the sensor API’s \code{sample_fetch} and \code{channel_get}.
  \item Zephyr modules: keeping drivers out of tree, with \code{zephyr/module.yml}.
  \item Emulators: a chip in software, on \code{native_sim}’s emulated I2C bus, and the sensor emulator
    backend.
  \item ztest: suites, tests, assertions, and hooks that run around every test.
  \item Twister: \code{testcase.yaml}, scenarios, and running a test suite in one command.
\end{itemize}
\end{chapterbox}

Zephyr's device model is what turns a node in the devicetree into a \code{struct device} the
application can use without knowing what is behind it. This chapter writes one: a binding, a Kconfig
option that turns itself on when the devicetree has the part, a driver that serves every instance of
it, and an implementation of the sensor API, so the application never names the chip at all.

Then it does what no board lets you do. The part is emulated, in software, on \code{native_sim}'s
I2C bus, and the emulator can be told to answer wrongly: a bad identity register, a conversion that
never finishes. That is what makes the driver testable, and \secref{c5:app:b} is about the testing:
ztest for the tests themselves, Twister for running them. The worked example is
\file{lectures/L05/examples/thermo}.

% Section 5.1, from lectures/L05/appendix/a_device_drivers.md.

\section{Device Drivers}
\label{c5:sec:devicedrivers}\label{c5:app:a}
This section covers how Zephyr turns a devicetree node into a device, what a driver provides for it,
and how a driver for a new part is packaged: binding, Kconfig option, and module. By the end you
should be able to explain every file in \file{lectures/L05/examples/thermo/module}, and why each one
is needed.

\subsection{What a driver is, in Zephyr}
Chapter~\ref{c2:ch} used devices without asking where they came from:
\code{DEVICE_DT_GET(DT_ALIAS(led0))} returned a \code{const struct device *}, and the GPIO API did
the rest. This chapter is the other side of that call.

A Zephyr driver is three things:

\begin{enumerate}
  \item \textbf{A set of functions} that implement one of Zephyr’s device APIs, such as GPIO, I2C or sensor,
    for one kind of part.
  \item \textbf{An init function}, which the kernel calls once at boot to bring an instance of the part up.
  \item \textbf{A macro expansion per devicetree node}, which creates a \code{struct device} for every node
    that the driver serves, at compile time, with the node’s properties already filled in.
\end{enumerate}

There is no registration at run time and no probing. By the time the program links, every device
exists as a constant in flash, and the only thing left to happen at boot is the init function.

\subsection{\code{struct device}}
Every device is a \code{struct device}, defined in \code{<zephyr/device.h>}. Its first five members
are the ones a driver writer deals with:

\begin{ccode}
struct device {
	const char *name;           /* The devicetree node's name, e.g. "thermo@48". */
	const void *config;         /* Per-instance config: const, from devicetree. */
	const void *api;            /* The function table: the driver's API struct. */
	struct device_state *state; /* Owned by the kernel: init's state and result. */
	void *data;                 /* Per-instance run-time data: mutable. */
	/* ... */
};
\end{ccode}

\noindent \textbf{\code{config} and \code{data} are the driver’s own types}, which the kernel stores
as \code{void *}. The difference between them matters on a microcontroller:

\begin{itemize}
  \item \textbf{\code{config} is \code{const}.} It holds what the devicetree says about the instance: its
    bus, its address, its properties. It is built at compile time and lives in flash.
  \item \textbf{\code{data} is not.} It holds what changes at run time: the last reading, a mutex, a state
    machine. It lives in RAM.
\end{itemize}

In \file{lectures/L05/examples/thermo/module/drivers/sensor/thermo/thermo.c} they are:

\begin{ccode}
struct thermo_config {
	struct i2c_dt_spec i2c;
	uint32_t conversion_timeout_ms;
};

struct thermo_data {
	int16_t raw;
};
\end{ccode}

\noindent and every driver function starts by getting them back out of the device:

\begin{ccode}
const struct thermo_config *config = dev->config;
struct thermo_data *data = dev->data;
\end{ccode}

\noindent \textbf{The name is the node’s name}, \code{thermo@48}, not its label. That is what
\code{dev->name} prints and what the logs show.

\subsection{The API struct}
Each kind of device has an API: a struct of function pointers declared by its subsystem. For the
sensor subsystem it is \code{struct sensor_driver_api} in \code{<zephyr/drivers/sensor.h>}, whose
two mandatory members are:

\begin{ccode}
sensor_sample_fetch_t sample_fetch;
sensor_channel_get_t channel_get;
\end{ccode}

\noindent The functions an application calls, such as \code{sensor_sample_fetch(dev)}, are thin
wrappers that cast \code{dev->api} and call through it:

\begin{ccode}
static inline int z_impl_sensor_sample_fetch(const struct device *dev)
{
	const struct sensor_driver_api *api =
		(const struct sensor_driver_api *)dev->api;

	return api->sample_fetch(dev, SENSOR_CHAN_ALL);
}
\end{ccode}

A driver fills in its own table with \code{DEVICE_API()}:

\begin{ccode}
static DEVICE_API(sensor, thermo_api) = {
	.sample_fetch = thermo_sample_fetch,
	.channel_get = thermo_channel_get,
};
\end{ccode}

\noindent \code{DEVICE_API(sensor, name)} declares a \code{const struct sensor_driver_api}, and also
places it in a linker section of its own, one per API class. That lets \code{DEVICE_API_IS(sensor,
dev)} check at run time whether a device really has a sensor API, rather than trusting whoever
passed the pointer.

This is Zephyr’s polymorphism: the application’s \code{sensor_sample_fetch()} works the same for
Thermo, for a Bosch accelerometer, or for any of the hundreds of sensor drivers in the tree, and
none of them knows about the others.

\subsection{From devicetree node to device}

\subsubsection{\code{DT_DRV_COMPAT} and instances}
A driver serves every node with a given \code{compatible}, and says which one at the top of its
file:

\begin{ccode}
#define DT_DRV_COMPAT qacademy_thermo
\end{ccode}

\noindent The value is the compatible \code{"qacademy,thermo"} in lowercase, with every character
that cannot be in a C identifier turned into an underscore. From then on, the \code{DT_INST_...}
macros count the \textbf{instances} of that compatible whose \code{status} is \code{okay}, numbered
0, 1, 2, … in an order the build decides. \code{DT_INST_PROP(inst, conversion_timeout_ms)} is
\code{DT_PROP()} for instance \code{inst}.

\subsubsection{One device per instance}
The driver writes a macro that defines everything for one instance, and expands it for all of them:

\begin{ccode}
#define THERMO_DEFINE(inst)                                                 \
    static struct thermo_data thermo_data_##inst;                           \
                                                                            \
    static const struct thermo_config thermo_config_##inst = {              \
        .i2c = I2C_DT_SPEC_INST_GET(inst),                                  \
        .conversion_timeout_ms = DT_INST_PROP(inst, conversion_timeout_ms), \
    };                                                                      \
                                                                            \
    SENSOR_DEVICE_DT_INST_DEFINE(inst, thermo_init, NULL,                   \
                                 &thermo_data_##inst,                       \
                                 &thermo_config_##inst, POST_KERNEL,        \
                                 CONFIG_SENSOR_INIT_PRIORITY, &thermo_api);

DT_INST_FOREACH_STATUS_OKAY(THERMO_DEFINE)
\end{ccode}

\noindent With two Thermo nodes in the devicetree, that is two \code{thermo_data}, two
\code{thermo_config}, and two devices; with none, it is nothing at all.
\code{I2C_DT_SPEC_INST_GET(inst)} fills in the bus device and the address from the node’s parent and
its \code{reg} property.

The device macro takes, in order:

\begin{booktable}{@{}LL@{}}
  \thead{Argument} & \thead{Meaning} \\
  \midrule
  \code{inst} & The instance. The macro turns it into the node, and the node into the device’s name and symbol. \\
  \code{thermo_init} & The init function. \\
  \code{NULL} & Power management actions, which Thermo has none of. \\
  \code{&thermo_data_\#\#inst} & \code{dev->data}. May be \code{NULL}. \\
  \code{&thermo_config_\#\#inst} & \code{dev->config}. \\
  \code{POST_KERNEL} & The init level. \\
  \code{CONFIG_SENSOR_INIT_PRIORITY} & The priority within the level. \\
  \code{&thermo_api} & \code{dev->api}. \\
\end{booktable}

\code{SENSOR_DEVICE_DT_INST_DEFINE()} is \code{DEVICE_DT_INST_DEFINE()} plus the sensor subsystem’s
bookkeeping; other subsystems’ drivers use \code{DEVICE_DT_INST_DEFINE()} directly.

\subsubsection{The symbol behind \code{DEVICE_DT_GET}}
A device is a global variable named after the node’s position in the devicetree’s dependency order:
\code{__device_dts_ord_32} is the device of the node with ordinal 32, and the ordinals are listed at
the top of \code{build/zephyr/include/generated/zephyr/devicetree_generated.h}.
\code{DEVICE_DT_GET(node)} is a reference to that variable.

That explains the error Chapter~\ref{c2:ch} met, from the other side. \code{DEVICE_DT_GET()} of an
enabled node that no driver defined a device for compiles, since the name is only a symbol, but the
link fails:

\begin{console}
undefined reference to `__device_dts_ord_32'
\end{console}

\noindent The node exists; nothing defined its device. The causes are a driver that is not enabled,
a driver whose \code{DT_DRV_COMPAT} does not match the node, or a module that is not in the build.
In the Thermo example, removing the module from \code{CMakeLists.txt} does exactly that.

\subsection{Initialization}

\subsubsection{Levels and priorities}
The kernel initializes devices in levels, in this order:

\begin{booktable}{@{}LL@{}}
  \thead{Level} & \thead{When} \\
  \midrule
  \code{EARLY} & Right after entering C, before the kernel’s architecture-specific initialization. \\
  \code{PRE_KERNEL_1}, \code{PRE_KERNEL_2} & Before the kernel’s services are available: no sleeping, no kernel objects. Interrupt controllers, clocks, the system timer, consoles. \\
  \code{POST_KERNEL} & After the kernel is running. Most device drivers, since they may need to sleep or take a mutex. \\
  \code{APPLICATION} & Just before \code{main()}. \\
\end{booktable}

Within a level, devices run in order of their \textbf{priority}, a number from 0 to 999, lowest
first. Subsystems give their drivers a Kconfig option for it: I2C controllers use
\code{CONFIG_I2C_INIT_PRIORITY}, which defaults to \code{CONFIG_KERNEL_INIT_PRIORITY_DEVICE} (50),
and sensors use \code{CONFIG_SENSOR_INIT_PRIORITY} (90). A sensor on an I2C bus therefore
initializes after the bus, which it needs.

\textbf{The build checks that.} \code{CONFIG_CHECK_INIT_PRIORITIES}, on by default, compares the
init order in the image with the dependencies in the devicetree, and fails the build if a device
initializes before a device its node depends on, such as a sensor before its bus.

To see the order, list the build’s init entries:

\begin{shell}
west build -t initlevels
\end{shell}

\noindent For the Chapter~\ref{c5:ch} exercises’ test application with two Lumo parts, the end of
the list is:

\begin{console}
POST_KERNEL
  __init_enable_logger: enable_logger(NULL)
  __init___device_dts_ord_31: i2c_emul_init(__device_dts_ord_31)
  __init___device_dts_ord_32: lumo_init(__device_dts_ord_32)
  __init___device_dts_ord_33: lumo_init(__device_dts_ord_33)
\end{console}

\noindent The list is read from the image, so it also shows a device whose initialization is
deferred (below), which the kernel then skips at boot.

\subsubsection{The init function}
An init function takes the device and returns 0 or a negative error number:

\begin{ccode}
static int thermo_init(const struct device *dev)
{
	const struct thermo_config *config = dev->config;
	uint8_t id;
	int ret;

	if (!i2c_is_ready_dt(&config->i2c)) {
		return -ENODEV;
	}

	ret = i2c_reg_read_byte_dt(&config->i2c, THERMO_REG_WHO_AM_I, &id);
	if (ret < 0) {
		return ret;
	}

	if (id != THERMO_WHO_AM_I_VALUE) {
		return -ENODEV;
	}

	return 0;
}
\end{ccode}

\noindent What belongs in it: checking that the bus is ready, that the part answers, and that it is
the part the devicetree says it is; and putting the part into the state the devicetree asks for.
What does not: anything slow that a driver could do on first use instead. Every init function runs
before \code{main()}, one after another.

\subsubsection{\code{device_is_ready()}}
The kernel records whether a device’s init function ran, and what it returned.
\code{device_is_ready(dev)} is true only if both went well:

\begin{ccode}
return dev->state->initialized && (dev->state->init_res == 0U);
\end{ccode}

\noindent \textbf{A failed init is final.} Nothing retries it, and every call into the device’s API
would talk to a part that is not there or not right. That is why every user of a device checks
\code{device_is_ready()} once, before using it, and why a driver’s init must fail when the part is
wrong rather than log a warning and carry on.

\subsubsection{Deferred initialization}
A node with the property \code{zephyr,deferred-init} is not initialized at boot. The application
initializes it later, when it chooses, with \code{device_init(dev)}: for example once a power rail
has been switched on. \code{device_init()} returns what the init function returned, or
\code{-EALREADY} if the device was initialized before, \textbf{whether or not that went well}. The
Thermo tests use exactly this to test a failing init: they break the emulated part, call
\code{device_init()}, and get \code{-ENODEV}; a second call, after repairing the part, gets
\code{-EALREADY}, and the device stays not ready.

\subsection{A binding and a Kconfig option for a new part}

\subsubsection{The binding}
Chapter~\ref{c2:ch} wrote a binding for a node of the application’s own. A driver’s binding is the
same kind of file, and says what a node for the part must look like:

\begin{yamlcode}
description: |
  Thermo, a made-up I2C temperature sensor ...

compatible: "qacademy,thermo"

include: [sensor-device.yaml, i2c-device.yaml]

properties:
  conversion-timeout-ms:
    type: int
    default: 100
\end{yamlcode}

\noindent \code{i2c-device.yaml} is Zephyr’s common binding for a part on an I2C bus: it requires
\code{reg}, the address, and makes the node valid only as a child of an I2C controller.
\code{sensor-device.yaml} adds the properties every sensor may have. A property with an \code{enum:}
list may only take one of the listed values, and \code{DT_INST_ENUM_IDX(inst, prop)} gives the index
of the value the node chose, at compile time.

\subsubsection{The vendor prefix}
The part before the comma in a compatible is a \textbf{vendor prefix}, and Zephyr keeps a list of
known ones in \code{dts/bindings/vendor-prefixes.txt}. A module can add its own, in a file of the
same name:

\begin{console}
qacademy	QRTECH Academy
\end{console}

\noindent Without it, a normal build only warns, once per node:

\begin{console}
node '/i2c@100/lumo@29' compatible 'qacademy,lumo' has unknown vendor prefix 'qacademy'
\end{console}

\noindent \textbf{Twister turns that warning into an error.} It configures every build with
\code{-DEXTRA_GEN_EDT_ARGS=--edtlib-Werror}, and the same message becomes a \code{devicetree error:}
that stops the configuration. A module that builds fine by hand and fails under Twister has usually
missed exactly this file.

\subsubsection{The Kconfig option}
Every driver has a Kconfig option that decides whether it is compiled:

\begin{kconfigcode}
config THERMO
	bool "Thermo I2C temperature sensor"
	default y
	depends on SENSOR
	depends on DT_HAS_QACADEMY_THERMO_ENABLED
	select I2C
\end{kconfigcode}

\noindent \textbf{\code{DT_HAS_QACADEMY_THERMO_ENABLED} is generated from the devicetree.} For every
compatible that has a binding, the build writes a symbol into \code{build/Kconfig/Kconfig.dts} that
is \code{y} when at least one node with that compatible is enabled:

\begin{console}
DT_COMPAT_QACADEMY_LUMO := qacademy,lumo

config DT_HAS_QACADEMY_LUMO_ENABLED
	def_bool $(dt_compat_enabled,$(DT_COMPAT_QACADEMY_LUMO))
\end{console}

\noindent With \code{default y}, the driver is on exactly when the hardware is there, and nobody
sets \code{CONFIG_THERMO} in a \code{prj.conf}. This is the pattern of nearly every driver in
Zephyr, and the reason enabling a driver usually means adding a devicetree node rather than setting
an option. The remaining dependencies are the subsystem the driver plugs into (\code{SENSOR}, which
the application enables) and the bus it needs (\code{I2C}, which the driver selects).

One consequence is easy to miss. A \code{default} that does not take is not an assignment, so
nothing warns about it: an application with a Thermo node but without \code{CONFIG_SENSOR=y} builds
its devicetree, leaves \code{THERMO} off, and fails at link time with \code{undefined reference to
'__device_dts_ord_32'}.

\subsection{The sensor API}

\subsubsection{Fetch and get}
Reading a sensor takes two calls:

\begin{ccode}
ret = sensor_sample_fetch(dev);
ret = sensor_channel_get(dev, SENSOR_CHAN_AMBIENT_TEMP, &value);
\end{ccode}

\noindent \textbf{\code{sample_fetch} talks to the part}: it starts a measurement if the part needs
one, reads the raw result, and stores it in \code{dev->data}. It can fail, and it can take time.
\code{sensor_sample_fetch(dev)} fetches every channel the part has (\code{SENSOR_CHAN_ALL}), and
\code{sensor_sample_fetch_chan(dev, chan)} only one.

\textbf{\code{channel_get} only converts} the stored result into the channel’s unit. It does no I/O,
and it returns the same value every time until the next fetch. The split lets a part that measures
several things at once, such as an accelerometer’s three axes, be read in one bus transaction and
then converted channel by channel. Both return \code{-ENOTSUP} for a channel the part does not have.

\subsubsection{\code{struct sensor_value}}
A reading is a fixed-point number in two parts:

\begin{ccode}
struct sensor_value {
	int32_t val1;   /* Integer part. */
	int32_t val2;   /* Fractional part, in millionths. */
};
\end{ccode}

\noindent Each channel has a unit: degrees Celsius for \code{SENSOR_CHAN_AMBIENT_TEMP}, lux for
\code{SENSOR_CHAN_LIGHT}. \textbf{Both parts carry the sign}: -1.5 degrees is \code{val1 = -1, val2
= -500000}, and -0.5 degrees is \code{val1 = 0, val2 = -500000}. A driver should not assemble the
two parts by hand;
\code{sensor_value_from_micro(&value, micro)} does it from a value in millionths, and
\code{sensor_value_to_micro(&value)} converts back. Work in \code{int64_t} millionths: a 16-bit count times
250000 millionths per count does not fit in 32 bits.

\subsubsection{The newer, asynchronous API}
Zephyr also has an asynchronous sensor API (\code{CONFIG_SENSOR_ASYNC_API}), built on RTIO, in which
a driver provides \code{submit} and \code{get_decoder} instead, and readings are decoded into Q31
fixed-point values. It suits high-rate and streaming sensors. The fetch and get API is what most
existing drivers implement, and what this book uses; Q31 values reappear in \secref{c5:app:b}, in
the emulator backend.

\subsection{Zephyr modules}

\subsubsection{Out of tree}
A driver for a part that Zephyr does not support does not have to live in Zephyr. A \textbf{module}
is a directory that the build treats as an extension of Zephyr: its CMake is included into Zephyr’s
build, its Kconfig into Zephyr’s configuration, and its bindings into the devicetree search path.
Vendor HALs, mbedTLS and LVGL are all modules; so is \code{examples/thermo/module}.

What makes a directory a module is one file, \code{zephyr/module.yml}:

\begin{yamlcode}
name: qacademy-thermo
build:
  cmake: .
  kconfig: Kconfig
  settings:
    dts_root: .
\end{yamlcode}

\begin{booktable}{@{}lL@{}}
  \thead{Key} & \thead{Meaning} \\
  \midrule
  \code{name} & The module’s name. The build also defines a Kconfig symbol for it: \code{CONFIG_ZEPHYR_QACADEMY_THERMO_MODULE=y}. \\
  \code{build: cmake} & The directory whose \code{CMakeLists.txt} the build includes. \\
  \code{build: kconfig} & The Kconfig file the build sources. \\
  \code{settings: dts_root} & A directory whose \code{dts/bindings/} holds bindings, and whose \code{dts/bindings/vendor-prefixes.txt} holds vendor prefixes. \\
\end{booktable}

The module’s \code{CMakeLists.txt} is included after Kconfig has run, so it compiles the driver only
when the driver is enabled:

\begin{cmakecode}
add_subdirectory_ifdef(CONFIG_THERMO drivers/sensor/thermo)
\end{cmakecode}

\noindent and the driver’s own \code{CMakeLists.txt} makes it a library of the Zephyr image, as
in-tree drivers do:

\begin{cmakecode}
zephyr_library()
zephyr_library_sources(thermo.c)
\end{cmakecode}

\subsubsection{Adding a module to a build}
In a west workspace, the modules are the projects of the manifest, as \secref{c1:app:b} described. A
module that is not a west project is added by path, in \code{ZEPHYR_EXTRA_MODULES}, before
\code{find_package(Zephyr)}:

\begin{cmakecode}
list(APPEND ZEPHYR_EXTRA_MODULES ${CMAKE_CURRENT_SOURCE_DIR}/../module)
find_package(Zephyr REQUIRED HINTS $ENV{ZEPHYR_BASE})
\end{cmakecode}

\noindent or on the command line, as \code{-DZEPHYR_EXTRA_MODULES=<path>}. The build merges the two,
so an application can name the modules it always needs in its \code{CMakeLists.txt}, and a test can
add another, such as an emulator, from outside.

\subsection{Summary}
\begin{itemize}
  \item A driver is an API implementation, an init function, and a macro expanded once for every enabled
    devicetree node of its compatible.
  \item A device’s \code{config} is \code{const}, from devicetree, in flash; its \code{data} changes at run
    time, in RAM.
  \item \code{DT_DRV_COMPAT} names the compatible; \code{DT_INST_FOREACH_STATUS_OKAY()} and
    \code{DEVICE_DT_INST_DEFINE()} make one device per instance.
  \item \code{undefined reference to '__device_dts_ord_N'} means a node that no driver made a device for.
  \item Devices initialize by level, then by priority. \code{west build -t initlevels} shows the order, and
    the build checks it against the devicetree.
  \item A failed init is final: \code{device_is_ready()} is false from then on. \code{zephyr,deferred-init}
    and \code{device_init()} move init to a moment the application chooses.
  \item A new part needs a binding, a vendor prefix, and a Kconfig option that depends on
    \code{DT_HAS_<COMPAT>_ENABLED}, so that the devicetree node alone enables the driver.
  \item The sensor API splits I/O (\code{sample_fetch}) from conversion (\code{channel_get}), and returns
    fixed-point \code{sensor_value}s whose two parts both carry the sign.
  \item A module packages drivers, bindings and Kconfig outside Zephyr, and is added to a build with
    \code{ZEPHYR_EXTRA_MODULES} or a west manifest.
\end{itemize}

% Section 5.2, from lectures/L05/appendix/b_emulators_and_testing.md.

\section{Emulators and Testing}
\label{c5:sec:emulatorsandtestin}\label{c5:app:b}
A driver for a part nobody has is untestable, unless something plays the part. This section covers
Zephyr’s emulators, which do, and the two tools that run tests against them: the ztest framework,
which the tests are written in, and Twister, which builds and runs them. By the end you should be
able to explain every file in
\file{lectures/L05/examples/thermo/module/drivers/sensor/thermo/thermo_emul.c} and
\file{lectures/L05/examples/thermo/tests}.

\subsection{Emulators}

\subsubsection{What an emulator is}
An \textbf{emulator} is a piece of software that stands in for a peripheral chip, at the level of
the bus the chip is on. The driver under test is the real driver, unchanged: it calls
\code{i2c_reg_read_byte_dt()} as it would on a board. What changes is below it. On
\code{native_sim}, the I2C controller \code{i2c0} is \code{zephyr,i2c-emul-controller}, a driver
that, instead of clocking bits onto wires, hands every transfer to the emulator registered at the
transfer’s address.

\begin{console}
application ── sensor API ──> thermo.c ── I2C API ──> i2c_emul ──> thermo_emul.c
                              (driver)                (i2c0)       (the "chip")
\end{console}

That is what makes emulators more useful than mocks: the driver’s I2C calls, its register addresses,
its byte order and the order of its transfers all run for real, and only the silicon is software.

The emulated I2C controller needs \code{CONFIG_EMUL=y}. Without it, \code{native_sim}'s \code{i2c0}
has no driver, and a sensor on that bus fails to link: in the exercises' application, with
\code{undefined reference to '__device_dts_ord_31'}, where node 31 is \code{i2c@100}.

\subsubsection{Defining one}
An emulator is defined per devicetree node, like a device, and usually for the same compatible as
the driver:

\begin{ccode}
#define DT_DRV_COMPAT qacademy_thermo

#define THERMO_EMUL_DEFINE(inst)                                            \
    static struct thermo_emul_data thermo_emul_data_##inst;                 \
    EMUL_DT_INST_DEFINE(inst, thermo_emul_init, &thermo_emul_data_##inst,   \
                        NULL, &thermo_emul_i2c_api, &thermo_emul_sensor_api)

DT_INST_FOREACH_STATUS_OKAY(THERMO_EMUL_DEFINE)
\end{ccode}

\noindent \code{EMUL_DT_INST_DEFINE()} takes the instance, an init function, the emulator’s data and
configuration, the \textbf{bus API} it answers transfers with, and an optional \textbf{backend API}
that tests use to control it. The result is a \code{struct emul}, which tests get with
\code{EMUL_DT_GET(node)}, just as they get the device with \code{DEVICE_DT_GET(node)}.

When the emulated I2C controller initializes, it looks up an emulator for each of its child nodes
and registers it at the node’s address. That is logged at boot:

\begin{console}
<inf> emul: Registering 1 emulator(s) for i2c@100
<inf> i2c_emul_ctlr: Register emulator 'lumo@29' at I2C addr 29
\end{console}

A node without an emulator is logged as \code{Cannot find emulator for 'lumo@29'}, and every
transfer to its address fails with \code{-EIO}: the bus behaves as if nothing answered there. The
controller initializes before the drivers on its bus, so every emulator is in its reset state by the
time a driver’s init function reads from it.

\subsubsection{The bus API}
For I2C, the bus API is one function:

\begin{ccode}
static const struct i2c_emul_api thermo_emul_i2c_api = {
	.transfer = thermo_emul_transfer,
};
\end{ccode}

\noindent It receives the messages of one I2C transfer, exactly as the driver’s call produced them.
The common register helpers arrive in two shapes:

\begin{itemize}
  \item \code{i2c_reg_write_byte_dt()} is \textbf{one write message}: the register address, then the value.
  \item \code{i2c_reg_read_byte_dt()} and \code{i2c_burst_read_dt()} are \textbf{two messages}: a write of
    the register address, then a read of one or more bytes, which by convention continue at the
    following registers.
\end{itemize}

An emulator’s transfer function is therefore a small register map: it keeps the registers in RAM,
updates them on writes, and fills in read buffers from them, doing whatever the datasheet says
happens in between, such as clearing a one-shot bit when a conversion finishes. Everything a real
part would get wrong-footed by, such as a write to a read-only register, the emulator decides how to
handle; Thermo’s ignores it and logs a warning, as the part would ignore it.

\subsubsection{The backend API}
Tests need to set what the part measures. They could call functions of the emulator’s own, and
Thermo’s emulator does have a header of its own, \code{thermo_emul.h}, for faults no real part has
on purpose: a wrong \code{WHO_AM_I}, a conversion that never finishes.

For the measured value itself there is a generic interface, the \textbf{sensor emulator backend} in
\code{<zephyr/drivers/emul_sensor.h>}:

\begin{ccode}
static const struct emul_sensor_driver_api thermo_emul_sensor_api = {
	.set_channel = thermo_emul_set_channel,
	.get_sample_range = thermo_emul_get_sample_range,
};
\end{ccode}

\noindent and tests call it through \code{emul_sensor_backend_set_channel()}, which works the same
for every sensor emulator in Zephyr that implements it. The value is passed as a \textbf{Q31}
fixed-point fraction with a shift: the physical value is \code{value / 2^31 * 2^shift}, in the
channel’s unit. A shift of 8 covers plus or minus 256 degrees; a shift of 15, 32768 lux. Converting
a value in millionths:

\begin{ccode}
const int8_t shift = 8;
q31_t value = (q31_t)((micro_c * (1LL << (31 - shift))) / 1000000);

emul_sensor_backend_set_channel(
	emul, (struct sensor_chan_spec){SENSOR_CHAN_AMBIENT_TEMP, 0}, &value, shift);
\end{ccode}

\noindent The emulator converts the value into what the part would register, rounded as the part
would round it. The test then reads it back through the driver and checks the conversion end to end.
The emulated I2C bus also has a lower-level escape hatch, \code{mock_api} on \code{struct i2c_emul},
which lets a test intercept transfers itself, for faults such as a bus that returns an error.

\subsubsection{What an emulator can and cannot show}
An emulator shows whether a driver talks to the part the way the datasheet says: the right
registers, the right byte order, the right conversion, the right reaction to a wrong ID or a
timeout. It shows it deterministically, on every commit, without hardware.

It cannot show what the author of the emulator did not know or did not model: a real part that needs
30 ms before it answers after power-up, a measurement that changes between reading two registers,
electrical noise, or a datasheet that is wrong. An emulator is a second implementation of the
datasheet, and a driver that passes against it agrees with that implementation. The first run on a
board is still a test.

\subsection{ztest}

\subsubsection{A test application}
A ztest test is an ordinary Zephyr application, built for a board, whose \code{main()} is provided
by the framework. \code{prj.conf} enables it:

\begin{kconfigcode}
CONFIG_ZTEST=y
\end{kconfigcode}

\noindent and the tests are functions in suites:

\begin{ccode}
#include <zephyr/ztest.h>

ZTEST_SUITE(thermo, NULL, NULL, NULL, after, NULL);

ZTEST(thermo, test_zero)
{
	zassert_ok(set_temperature(0));

	struct sensor_value value = read_temperature();

	zassert_equal(value.val1, 0);
	zassert_equal(value.val2, 0);
}
\end{ccode}

\noindent On \code{native_sim}, the program runs every suite, prints a summary, and exits with
status 0 if every test passed and 1 if any failed, which is what Twister, and the repository’s
\code{make test}, check. Tests built with ztest also run in simulated time as fast as the host
allows: \code{native_sim} only slows down to real time by default when the build is not a test.

\begin{console}
SUITE PASS - 100.00% [thermo]: pass = 11, fail = 0, skip = 0, total = 11 duration = 0.000 seconds
 - PASS - [thermo.test_channel_get_returns_the_last_fetch] duration = 0.000 seconds
 ...
PROJECT EXECUTION SUCCESSFUL
\end{console}

\subsubsection{Suites and hooks}
\code{ZTEST_SUITE(name, predicate, setup, before, after, teardown)} declares a suite and,
optionally, functions that run around it:

\begin{booktable}{@{}lL@{}}
  \thead{Argument} & \thead{Runs} \\
  \midrule
  \code{predicate} & Before the suite, to decide whether it runs at all. \\
  \code{setup} & Once, before the suite’s first test. Its return value is the suite’s \emph{fixture}. \\
  \code{before} & Before every test. \\
  \code{after} & After every test, whether it passed or failed. \\
  \code{teardown} & Once, after the suite’s last test. \\
\end{booktable}

\code{ZTEST(suite, name)} defines a test; \code{ZTEST_F(suite, name)} defines one that receives the
fixture. Thermo’s suite uses \code{after} to repair whatever fault a test injected, so that a
failing test cannot make the next one fail too.

\textbf{Tests run in alphabetical order}, not in the order they appear in the file. A test that only
passes because another ran first is a test that will one day fail for no reason; every test must set
up everything it depends on, and the hooks exist for exactly that.

\subsubsection{Assertions}
\begin{booktable}{@{}LL@{}}
  \thead{Macro} & \thead{On failure} \\
  \midrule
  \code{zassert_true(cond, msg, ...)}, \code{zassert_false}, \code{zassert_equal(a, b, ...)}, \code{zassert_ok(expr)}, … & Fails the test and ends it. \\
  \code{zexpect_true}, \code{zexpect_equal}, … & Fails the test, but lets it continue, to report more than the first failure. \\
  \code{zassume_true}, \code{zassume_equal}, … & Skips the rest of the test: for a precondition that is not what is being tested. \\
\end{booktable}

Every assertion takes an optional \code{printf}-style message, which is printed with the file, the
line and the failed condition. A good message says what should have happened, in the terms of the
datasheet or the exercise, not only that two numbers differed:

\begin{console}
Assertion failed at .../conversion/src/main.c:71: conversion_test_saturated: (value.val1 not equal to 16383)
0xFFFF counts is 16383.75 lux, not -796.-119184
\end{console}

\subsubsection{What to test in a driver}
Against an emulator, a driver’s tests cover what a datasheet specifies and what a user relies on:

\begin{itemize}
  \item \textbf{Readiness}: the device is ready when the part is genuine, and not when it is not.
  \item \textbf{Conversion at the edges}: zero, the smallest step, negative values if the part has them, and
    the limits of the range, where overflows and sign errors live.
  \item \textbf{Configuration from devicetree}: what a property should put into a register, read back from
    the emulator.
  \item \textbf{The API contract}: unsupported channels, fetch without get and get without fetch.
  \item \textbf{Faults}: a timeout, a wrong ID, a bus error. These are the tests only an emulator makes
    cheap.
\end{itemize}

\subsection{Twister}

\subsubsection{Test applications as scenarios}
\textbf{Twister} is Zephyr’s test runner, and the same tool the repository’s \code{make examples}
uses. It walks a directory tree, finds every application with a \code{testcase.yaml} (or, for
samples, a \code{sample.yaml}), and builds, runs and checks each of its \textbf{scenarios} on each
selected platform, in parallel.

The Thermo tests’ \code{testcase.yaml}:

\begin{yamlcode}
common:
  platform_allow: native_sim/native/64
  integration_platforms:
    - native_sim/native/64
  harness: ztest
  tags:
    - sensor
tests:
  qacademy.l05.thermo.driver: {}
\end{yamlcode}

\begin{booktable}{@{}LL@{}}
  \thead{Key} & \thead{Meaning} \\
  \midrule
  \code{tests} & The scenarios, by name. Each builds the application once; keys under a scenario override \code{common}. \\
  \code{platform_allow} & The platforms the scenario may run on. \\
  \code{integration_platforms} & The platforms it runs on when Twister is run with \code{--integration}, as Zephyr’s CI does for pull requests. \\
  \code{harness: ztest} & How to judge the run: parse ztest’s output, and report every test case. A sample uses \code{harness: console} with regular expressions instead. \\
  \code{extra_configs}, \code{extra_args} & Kconfig settings and CMake arguments for the scenario, such as \code{CONFIG_ASSERT=y} or \code{EXTRA_CONF_FILE=...}. \\
  \code{tags}, \code{filter} & Selecting scenarios: by tag on the command line, or by an expression over the build’s configuration. \\
\end{booktable}

A suite and a scenario are different things. A ztest \textbf{suite} is a group of tests inside one
program. A Twister \textbf{scenario} is one build of the program: the same tests built twice, once
with \code{CONFIG_FOO=y} and once without, are two scenarios running the same suites.

\subsubsection{Running it}
\begin{shell}
west twister -T lectures/L05/examples/thermo/tests -p native_sim/native/64
\end{shell}

\code{-T} gives the directory to search, and \code{-p} the platform. Twister writes everything to
\code{twister-out/} in the current directory, renaming an older one to \code{twister-out.1}: one
build directory per scenario, under
\code{twister-out/native_sim_native_64/host_gnu/<path>/<scenario>/}, with the program’s output in
\code{handler.log} and the build’s in \code{build.log}, and reports in
\code{twister-out/twister.json} and \code{twister.xml}. \code{-i} (\code{--inline-logs}) prints the
log of every failure in the terminal instead.

\textbf{Twister is stricter than a plain build.} It builds with compiler warnings as errors
(\code{CONFIG_COMPILER_WARNINGS_AS_ERRORS=y}) and with devicetree warnings as errors
(\code{--edtlib-Werror}), which is how a missing \code{vendor-prefixes.txt} that only warned by hand
stops a Twister run. \code{--disable-warnings-as-errors} turns both off, and should not be needed.

Extra CMake arguments for every build go after \code{-x}, for example a module that is not in the
application’s \code{CMakeLists.txt}:

\begin{shell}
west twister -T <tests> -p native_sim/native/64 \
    -x=ZEPHYR_EXTRA_MODULES=<path to module>
\end{shell}

\noindent With \code{--coverage}, Twister builds with gcov instrumentation and collects coverage of
the code the tests ran into an HTML report. Chapter~\ref{c6:ch} shows it.

The same \code{testcase.yaml} runs on real boards too, with \code{--device-testing} and a serial
port: the scenarios that pass here against an emulator are what \textbf{Hardware-in-the-Loop
Testing} runs against the real part.

\subsection{Summary}
\begin{itemize}
  \item An emulator plays a peripheral chip at its bus. The driver under test is unchanged; on
    \code{native_sim}, the emulated I2C controller hands each transfer to the emulator at its address.
  \item \code{EMUL_DT_INST_DEFINE()} defines an emulator per devicetree node, with a bus API that answers
    transfers and an optional backend API that tests control it through. \code{EMUL_DT_GET()} gets it.
  \item The sensor emulator backend sets a measured value as Q31 with a shift, the same way for every sensor
    emulator. Faults no real part has on purpose go in the emulator’s own header.
  \item An emulator is a second implementation of the datasheet: it tests the driver’s agreement with it,
    not with the silicon.
  \item ztest applications are Zephyr applications; \code{ZTEST_SUITE()} declares a suite and its hooks,
    \code{ZTEST()} a test. Tests run in alphabetical order, so none may depend on another.
  \item \code{zassert_*} ends a test, \code{zexpect_*} lets it continue, and \code{zassume_*} skips it. Say
    in the message what should have happened.
  \item Twister finds \code{testcase.yaml} files, builds every scenario, runs it, and judges it by its
    harness, with warnings as errors. A suite is tests in a program; a scenario is a build of a program.
\end{itemize}

% Section 5.3, from lectures/L05/README.md: the objectives, the evaluation
% questions, and the next lecture.

\section{Review}
\label{c5:sec:review}

\needspace{6\baselineskip}
\subsection*{What you should be able to do now}
After this chapter, you should:

\begin{itemize}
  \item Understand how Zephyr turns a devicetree node into a \code{struct device} at compile time.
  \item Be able to write a driver that serves every instance of its compatible in the devicetree.
  \item Understand the order devices are initialized in, and how to find it for a build.
  \item Be able to implement a subsystem API, so that applications use the part without knowing it.
  \item Be able to package drivers, bindings and Kconfig as a Zephyr module.
  \item Be able to write an emulator that lets a driver run, and fail, without hardware.
  \item Be able to test a driver with ztest, and run the tests with Twister.
\end{itemize}

\needspace{6\baselineskip}
\subsection*{Questions to test yourself}
You should be able to explain:

\begin{enumerate}
  \item What is the difference between a device’s \code{config} and its \code{data}, and where does each one
    live?
  \item How does \code{DT_DRV_COMPAT} connect a driver to the devicetree nodes it serves?
  \item What decides the order in which devices are initialized, and how can you see it for a build?
  \item What does \code{device_is_ready()} return for a device whose init function failed?
  \item Why does the driver’s Kconfig option depend on \code{DT_HAS_QACADEMY_THERMO_ENABLED}?
  \item Why does a sensor driver split reading into \code{sample_fetch} and \code{channel_get}?
  \item What does an emulator make testable that a real part does not?
  \item What is the difference between a ztest suite and a Twister scenario?
  \item What goes wrong without the module’s \code{vendor-prefixes.txt}, and why only under Twister?
\end{enumerate}

\needspace{6\baselineskip}
\subsection*{Looking ahead}
Chapter~\ref{c6:ch} is about subsystems and debugging:

\begin{itemize}
  \item The shell, with commands of your own.
  \item Settings that survive a restart, on the simulated flash.
  \item Networking through the host’s sockets.
  \item Debugging on the host: GDB, sanitizers, fatal errors and coverage.
\end{itemize}

You will make the program talk back, remember things, and crash usefully.

% Section 5.4, from lectures/L05/appendix/c_exercises.md.

\section{Exercises}
\label{c5:sec:exercises}\label{c5:app:c}
These exercises reinforce the concepts from \secref{c5:app:a} and \secref{c5:app:b}.

\begin{aside}
\textbf{How to check your work.} Write each exercise set in its own directory under
\file{lectures/L05/exercises}, with the file names given there, and run \code{make test} from the
repository root. A set’s tests switch on as soon as its files exist, and a set you have not started
is reported as \code{SKIP}. Sets 1 to 3 are one driver growing: start each by copying the previous
set’s directory. The tests of Sets 1 to 3 are ztest applications in
\file{lectures/L05/exercises/test} that build your module against the course’s emulator; Set 4 is
checked by its output, and Set 5 by running your own tests.

The solutions are published in the companion repository, under
\file{lectures/L05/appendix/solutions}, a few days after the lecture. Try each exercise yourself
before you read them.
\end{aside}

Remember \code{source env.sh} in every new terminal.

\subsection{The part: Lumo}
Lumo is a made-up I2C ambient light sensor. It measures continuously: its DATA registers always hold
the latest measurement, and there is nothing to start. Its datasheet, in full:

\begin{booktable}{@{}lllL@{}}
  \thead{Register} & \thead{Name} & \thead{Access} & \thead{Contents} \\
  \midrule
  \code{0x00} & \code{CONFIG} & R/W & Bits 1:0, \code{GAIN}: \code{0} is a gain of 1, \code{1} of 2, \code{2} of 4, \code{3} of 8. Bits 7:2 are reserved, read as 0 and ignore writes. Reset value \code{0x00}. \\
  \code{0x01} & \code{DATA_LSB} & R & The measurement, low byte. \\
  \code{0x02} & \code{DATA_MSB} & R & The measurement, high byte. \\
  \code{0x0A} & \code{PART_ID} & R & Always \code{0x4C} on a genuine part. \\
\end{booktable}

\begin{itemize}
  \item The measurement is an \textbf{unsigned 16-bit count, little-endian}: \code{DATA_LSB} then
    \code{DATA_MSB}.
  \item At a gain of 1, \textbf{one count is 0.25 lux}. At a gain of G, one count is 0.25 / G lux.
  \item A count of \code{0xFFFF} means the part is \textbf{saturated}: the light is at least that bright.
  \item A read of several bytes continues at the following registers, so both DATA bytes can be read in one
    transfer. Writes to read-only registers are ignored.
\end{itemize}

The course provides the part itself: an emulator, in \file{lectures/L05/exercises/test/lumo_emul},
which is a Zephyr module of its own. It plays a Lumo chip for every enabled \code{"qacademy,lumo"}
node on \code{native_sim}’s I2C bus. The tests add it to their builds; to build with it by hand, add
it on the command line:

\begin{shell}
-DZEPHYR_EXTRA_MODULES=$PWD/lectures/L05/exercises/test/lumo_emul
\end{shell}

\noindent Tests control it through that module's own header, \code{<qacademy/lumo_emul.h>}
(\code{lumo_emul_set_part_id()}, \code{lumo_emul_set_counts()} and \code{lumo_emul_get_config()}),
and through the generic sensor emulator backend, with \code{SENSOR_CHAN_LIGHT} in lux.

% ---------------------------------------------------------------------
\exerciseset{A module and a device}
\label{c5:set:1}
Create a directory named \code{exercise1} to work in for this exercise set. The directory is your
module.

\exercise{The module and its Kconfig option}{Code}
\label{c5:ex:1.1}
In this exercise you will create a module that the build can include.

\task{Tasks}
Create these files in \code{exercise1}:

\begin{itemize}
  \item \code{zephyr/module.yml}: a module named \code{qacademy-lumo}, whose CMake and Kconfig are in the
    module’s top directory, and whose \code{dts_root} is the top directory too.
  \item \code{Kconfig}: an option \code{LUMO}, a \code{bool} with the prompt \code{Lumo I2C ambient light
    sensor}, which:
\begin{itemize}
  \item is \code{y} by default,
  \item depends on \code{SENSOR}, and on the devicetree having an enabled \code{"qacademy,lumo"} node,
  \item selects \code{I2C}.
\end{itemize}
  \item \code{CMakeLists.txt}: when \code{CONFIG_LUMO} is enabled, a Zephyr library holding
    \code{src/lumo.c}.
\end{itemize}

\exercise{The binding}{Code}
\label{c5:ex:1.2}

\task{Tasks}
\begin{itemize}
  \item Create \code{dts/bindings/qacademy,lumo.yaml}: a binding for the compatible \code{"qacademy,lumo"},
    which includes \code{sensor-device.yaml} and \code{i2c-device.yaml}, and has no properties of its
    own yet.
  \item Create \code{dts/bindings/vendor-prefixes.txt}, registering the vendor prefix \code{qacademy}.
\end{itemize}

\exercise{The driver’s init}{Code}
\label{c5:ex:1.3}

\task{Tasks}
\begin{itemize}
  \item Create \code{src/lumo_reg.h} with macros for the register addresses, the \code{GAIN} mask, and the
    value of \code{PART_ID}, from the datasheet above.
  \item Create \code{src/lumo.c}, the driver for the compatible \code{"qacademy,lumo"}:
\begin{itemize}
  \item A config struct holding the part’s \code{struct i2c_dt_spec}.
  \item A sensor API table, defined with \code{DEVICE_API(sensor, ...)}, whose \code{sample_fetch} and
    \code{channel_get} functions only return \code{-ENOTSUP} for now.
  \item An init function that:
\begin{itemize}
  \item returns \code{-ENODEV} if the I2C bus is not ready,
  \item reads \code{PART_ID}, and returns the error if the read fails,
  \item returns \code{-ENODEV} if \code{PART_ID} is not \code{0x4C},
  \item and otherwise returns 0.
\end{itemize}
  \item One device per enabled instance, defined with \code{SENSOR_DEVICE_DT_INST_DEFINE()} at
    \code{POST_KERNEL} and \code{CONFIG_SENSOR_INIT_PRIORITY}. The data pointer may be \code{NULL} for
    now.
\end{itemize}
\end{itemize}

The tests in \file{lectures/L05/exercises/test/device} put one part at address \code{0x29}, and a
second at \code{0x2a} whose initialization is deferred, so that a test can give it the wrong
\code{PART_ID} first. To build and run them by hand:

\begin{shell}
west build -p -b native_sim/native/64 -d build/device \
    lectures/L05/exercises/test/device -- \
    -DEXERCISE_DIR=$PWD/lectures/L05/exercises/exercise1
./build/device/zephyr/zephyr.exe
\end{shell}

\noindent The four tests must pass: the option turns itself on, the part is ready, its device has a
sensor API, and the wrong \code{PART_ID} makes \code{device_init()} return \code{-ENODEV}.

\task{Reflection}
\begin{itemize}
  \item Remove \code{vendor-prefixes.txt} and build the tests again by hand. What does the build print? Then
    build once more with \code{-DEXTRA_GEN_EDT_ARGS=--edtlib-Werror} added, as Twister does. What
    happens?
  \item Run \code{west build -d build/device -t initlevels}. Where do your two \code{lumo_init} calls come
    in the list, and what comes just before them? Why must it be in that order?
  \item The deferred part is in the \code{initlevels} list too. Why is it still not ready after boot?
  \item Find \code{DT_HAS_QACADEMY_LUMO_ENABLED} in the build directory. Which file defines it, and why
    could it not exist before you wrote the binding?
\end{itemize}

% ---------------------------------------------------------------------
\exerciseset{Reading light}
\label{c5:set:2}
Copy \code{exercise1} to a directory named \code{exercise2} to work in for this exercise set.

\exercise{Fetch and get}{Code}
\label{c5:ex:2.1}
In this exercise you will implement the sensor API, at the part’s reset gain of 1.

\task{Tasks}
In \code{src/lumo.c}:

\begin{itemize}
  \item Add a data struct holding the count of the last fetch, and pass it to the device.
  \item \code{sample_fetch} shall:
\begin{itemize}
  \item return \code{-ENOTSUP} for any channel other than \code{SENSOR_CHAN_ALL} and
    \code{SENSOR_CHAN_LIGHT},
  \item read both DATA bytes \textbf{in one transfer}, and return the error if the read fails,
  \item store the count, and return 0.
\end{itemize}
  \item \code{channel_get} shall:
\begin{itemize}
  \item return \code{-ENOTSUP} for any channel other than \code{SENSOR_CHAN_LIGHT},
  \item convert the stored count to lux, at 0.25 lux per count, into the \code{struct sensor_value}.
\end{itemize}
  \item Do no I/O in \code{channel_get}.
\end{itemize}

The tests in \file{lectures/L05/exercises/test/conversion} check, among others, that:

\begin{narrowtable}{@{}llll@{}}
  \thead{Count} & \thead{Lux} & \thead{\code{val1}} & \thead{\code{val2}} \\
  \midrule
  \code{0} & 0 & \code{0} & \code{0} \\
  \code{1} & 0.25 & \code{0} & \code{250000} \\
  \code{0x0102} & 64.5 & \code{64} & \code{500000} \\
  \code{0xFFFF} & 16383.75 & \code{16383} & \code{750000} \\
\end{narrowtable}

The tests of Set 1 must still pass.

\textbf{Hint:} \code{sys_get_le16()}, from \code{<zephyr/sys/byteorder.h>}, turns two little-endian
bytes into a \code{uint16_t}. \code{sensor_value_from_micro()} fills in a \code{struct sensor_value}
from millionths.

\task{Reflection}
\begin{itemize}
  \item Replace the single transfer with two \code{i2c_reg_read_byte_dt()} calls, one per DATA byte. Do the
    tests still pass? What could go wrong on a real part that the emulator cannot show?
  \item Compute the lux in \code{int} instead of \code{int64_t}, by removing the cast. Which test fails?
    Above how many counts does the multiplication overflow, and how bright is that?
  \item \code{test_channel_get_returns_the_last_fetch} changes the part’s count between a fetch and a get.
    Why is the old value the right answer?
\end{itemize}

% ---------------------------------------------------------------------
\exerciseset{The gain}
\label{c5:set:3}
Copy \code{exercise2} to a directory named \code{exercise3} to work in for this exercise set.

\exercise{A devicetree property}{Code}
\label{c5:ex:3.1}
In this exercise you will make the gain a property of the devicetree node, written to the part at
init.

\task{Tasks}
\begin{itemize}
  \item In the binding, add a property \code{gain}: an \code{int}, allowed to be only 1, 2, 4 or 8, with a
    default of 1.
  \item In \code{src/lumo.c}:
\begin{itemize}
  \item Add the gain and its \code{GAIN} bits to the config struct. Take the bits from the property’s index
    in the binding’s list of allowed values, at compile time.
  \item In the init function, after checking \code{PART_ID}, write the \code{GAIN} bits to \code{CONFIG},
    without changing the register’s other bits, and return the error if that fails.
  \item In \code{channel_get}, divide by the gain: at a gain of G, one count is 0.25 / G lux.
\end{itemize}
\end{itemize}

The tests in \file{lectures/L05/exercises/test/gain} put three parts on the bus, at gains of 4, 8
and the default, and check what each one’s \code{CONFIG} register holds after boot and what its
counts convert to: one count at a gain of 4 is 0.0625 lux, and \code{0xFFFF} at a gain of 8 is
2047.96875 lux. The tests of Sets 1 and 2 must still pass.

\textbf{Hint:} \code{DT_INST_ENUM_IDX(inst, prop)} is the index of a property’s value in its
binding’s \code{enum} list. \code{i2c_reg_update_byte_dt()} changes only the bits of a register that
a mask selects.

\task{Reflection}
\begin{itemize}
  \item Add \code{gain = <3>;} to a part in an overlay and build the gain tests. What happens, and at which
    step of the build?
  \item Why take the \code{GAIN} bits from \code{DT_INST_ENUM_IDX()}, rather than computing a base-2
    logarithm of the gain at run time?
  \item At a gain of 8, what is the smallest step between two readings, and the brightest light the part can
    measure? What are they at a gain of 1?
\end{itemize}

% ---------------------------------------------------------------------
\exerciseset{An application}
\label{c5:set:4}
Create an application named \code{exercise4} to work in for this exercise set.

\exercise{Reading once a second}{Code}
\label{c5:ex:4.1}

\task{Tasks}
\begin{itemize}
  \item \code{CMakeLists.txt}: add the module in \code{../exercise3} to the build, before finding Zephyr. Do
    \textbf{not} add the emulator.
  \item \code{prj.conf}: enable the sensor subsystem, emulators, and logging.
  \item \code{app.overlay}: a Lumo part labelled \code{light}, at address \code{0x29} on \code{&i2c0}, at the default gain.
  \item \code{src/main.c}, with a log module named \code{light_app} at level \code{INF}:
\begin{itemize}
  \item If the part is not ready, log \code{<name> is not ready} as an error, with the device’s name, and
    return.
  \item Log \code{Reading <name>}.
  \item Then, forever: fetch and get a light reading, log it as \code{<lux> lux} with two decimals, or
    \code{Reading failed: <error>} as an error, and sleep for one second.
\end{itemize}
\end{itemize}

Build it with the emulator, and with the emulator’s sunrise, which changes the light half a second
before each whole second:

\begin{shell}
west build -p -b native_sim/native/64 lectures/L05/exercises/exercise4 -- \
    -DZEPHYR_EXTRA_MODULES=$PWD/lectures/L05/exercises/test/lumo_emul \
    -DCONFIG_EMUL_LUMO_DAYLIGHT=y
./build/zephyr/zephyr.exe -stop_at=6.5
\end{shell}

\noindent The output must contain, in this order, with the emulator’s own log lines in between:

\begin{console}
[00:00:00.000,000] <inf> light_app: Reading lumo@29
[00:00:00.000,000] <inf> light_app: 0.00 lux
[00:00:01.010,000] <inf> light_app: 0.25 lux
[00:00:02.020,000] <inf> light_app: 42.00 lux
[00:00:03.030,000] <inf> light_app: 320.00 lux
[00:00:04.040,000] <inf> light_app: 2750.00 lux
[00:00:05.050,000] <inf> light_app: 12000.00 lux
[00:00:06.060,000] <inf> light_app: 16383.75 lux
\end{console}

\task{Reflection}
\begin{itemize}
  \item Build the application without \code{-DZEPHYR_EXTRA_MODULES=...}. Does it build? What does it print,
    and why?
  \item Build it with the emulator, but with \code{-DCONFIG_EMUL=n} added. What happens, and which
    devicetree node is the error about?
  \item The emulator’s last light level is 30000 lux, but the reading is 16383.75. Why? Set \code{gain =
    <8>} in \code{app.overlay} and run it again: what do the readings become, and what does that say
    about choosing a gain?
\end{itemize}

% ---------------------------------------------------------------------
\exerciseset{Tests of your own}
\label{c5:set:5}
Create a directory named \code{exercise5} to work in for this exercise set: a ztest application that
tests the driver of Exercise Set~\ref{c5:set:3}.

\exercise{A test suite and a Twister scenario}{Code}
\label{c5:ex:5.1}

\task{Tasks}
\begin{itemize}
  \item \code{CMakeLists.txt}: add the module in \code{../exercise3} to the build, and \code{src/main.c} to
    the application. As in Set 4, do not add the emulator.
  \item \code{prj.conf}: enable ztest, the sensor subsystem and emulators.
  \item \code{app.overlay}: two Lumo parts, one at the default gain and one at a gain of 2.
  \item \code{src/main.c}: a suite named \code{lumo}, whose tests check at least:
\begin{itemize}
  \item that both parts are ready,
  \item counts of 0, 1 and \code{0xFFFF} at the default gain,
  \item a count at a gain of 2,
  \item that the gain of 2 was written to the part’s \code{CONFIG} register,
  \item and that an unsupported channel returns \code{-ENOTSUP} from both fetch and get.
\end{itemize}
  \item \code{testcase.yaml}: one scenario, named \code{qacademy.l05.exercise5.lumo}, allowed only on
    \code{native_sim/native/64}, judged by the ztest harness.
\end{itemize}

Build and run it with the emulator:

\begin{shell}
west build -p -b native_sim/native/64 lectures/L05/exercises/exercise5 -- \
    -DZEPHYR_EXTRA_MODULES=$PWD/lectures/L05/exercises/test/lumo_emul
./build/zephyr/zephyr.exe
\end{shell}

\noindent The suite must pass: the output must contain \code{SUITE PASS - 100.00\% [lumo]} and
\code{PROJECT EXECUTION SUCCESSFUL}.

\task{Reflection}
\begin{itemize}
  \item Run the scenario with Twister:
\begin{shell}
west twister -T lectures/L05/exercises/exercise5 -p native_sim/native/64 \
    -x=ZEPHYR_EXTRA_MODULES=$PWD/lectures/L05/exercises/test/lumo_emul
\end{shell}
Where does Twister put the output of your test program? What does it report besides pass and fail?
  \item Break your Set 3 driver on purpose: make \code{channel_get} ignore the gain. Which of your tests
    fail? If none do, which test is missing?
  \item Your tests and the course’s tests both pass. What would you still check on the first board with a
    real Lumo part?
\end{itemize}

